\documentclass[a4paper]{article}

% Global layout
\usepackage{fancyhdr, graphicx, hyperref, indentfirst, lastpage, setspace}
\usepackage[margin=40mm]{geometry}

% Encoding
\usepackage[utf8]{vntex, inputenc}
\usepackage[english]{babel}
\usepackage{amsmath, amssymb, gensymb}

% Better table
\usepackage{array, booktabs, multicol, multirow, siunitx, tabularx}

% Code space
\usepackage[dvipsnames]{xcolor}
\usepackage{tikz}
\usepackage[framemethod=tikz]{mdframed}
\usepackage{minted, verbatim} % needs --shell-escape flag and Pygments

% Graphics
\usepackage{caption, float}

% Page setup
\allowdisplaybreaks{} % to have page breaks inside align* environment
\hypersetup{urlcolor=blue,linkcolor=black,citecolor=red,colorlinks=true}
\usemintedstyle{emacs}
\numberwithin{equation}{section}
\renewcommand{\arraystretch}{1.2} % space between table rows

% Global style setup
\makeatletter % change font size for not having underfull hbox
\renewcommand\Huge{\@setfontsize\Huge{22pt}{18}}
\makeatother

\AtBeginDocument{\renewcommand*\contentsname{Contents}}
\AtBeginDocument{\renewcommand*\refname{References}}
\setlength{\headheight}{40pt}
\pagestyle{fancy}
\fancyhead{} % clear all header fields
\fancyhead[L]{
  \begin{tabular}{rl}
    \begin{picture}(25,15)(0,0)
    \put(0,-8){\includegraphics[width=8mm, height=8mm]{hcmut.png}}
    \end{picture}
    \begin{tabular}{l}
      \textbf{\bf \ttfamily University of Technology, Ho Chi Minh City}\\
      \textbf{\bf \ttfamily Faculty of Computer Science and Engineering}
    \end{tabular}
  \end{tabular}
}
\fancyhead[R]{
	\begin{tabular}{l}
		\tiny \bf \\
		\tiny \bf
	\end{tabular}  }
\fancyfoot{} % clear all footer fields
\fancyfoot[L]{\scriptsize \ttfamily Report for Advanced Programming --- Academic year 2020--2021}
\fancyfoot[R]{\scriptsize \ttfamily Page {\thepage}/\pageref{LastPage}}
\renewcommand{\headrulewidth}{0.3pt}
\renewcommand{\footrulewidth}{0.3pt}

% \everymath{\color{blue}}

\newcommand*\mean[1]{\bar{#1}}
\newenvironment{code}[1]
{\VerbatimEnvironment%
  \begin{mdframed}[leftline=false,rightline=false,backgroundcolor=magenta!10,nobreak=false]%
    \begin{minted}[linenos=true,breaklines,breaksymbolleft=,obeytabs=true,tabsize=2]{#1}%
}
{
    \end{minted}%
  \end{mdframed}%
}

\begin{document}

\begin{titlepage}
  \begin{center}
    VIETNAM NATIONAL UNIVERSITY, HO CHI MINH CITY \\
    UNIVERSITY OF TECHNOLOGY \\
    FACULTY OF COMPUTER SCIENCE AND ENGINEERING
  \end{center}

  \vspace{1cm}

  \begin{figure}[H]
    \centering
    \includegraphics[width=0.5\textwidth]{hcmut.png}
  \end{figure}

  \vspace{1cm}

  \begin{center}
    \begin{tabular}{c}
      \textbf{\Large Advanced Programming (CO2039)}             \\
      {}                                                        \\
      \midrule                                                  \\
      \textbf{\Large Report (Semester 202, Duration: 01 weeks)} \\
      {}                                                        \\
      \textbf{\Huge OOP vs FP}                                  \\
      {}                                                        \\
      \bottomrule
    \end{tabular}
  \end{center}

  \vspace{3cm}

  \begin{table}[h]
    \begin{tabular}{rl}
      \hspace{1cm} Advisor:      & Mr.\ Lê Lam Sơn \\
                                 &                 \\
      \hspace{1cm} Student Name: & Nguyễn Hoàng    \\
      \hspace{1cm} Student ID\@: & 1952255         \\
    \end{tabular}
  \end{table}

  \begin{center}
    {\footnotesize HO CHI MINH CITY, AUGUST 2021}
  \end{center}
\end{titlepage}


% \thispagestyle{empty}

% \newpage
% \tableofcontents
% \newpage

\section{OOP and FP in baking a pizza}
OOP makes code understandable by encapsulating moving parts.
FP makes code understandable by minimizing moving parts.

What? Alright that sounds a bit rough, let's rephrase this a bit.
OOP aims to model the world in self-contained entities, and affects change by modifying the state of itself or other entities.
FP on the other hand aims to not modify the original data, but rather creates new data given some existing data.

To demonstrate this, we will try to make a pizza.
With OOP, a big box or object with all the materials steps to make a pizza is available.
FP takes a different approach, as it only contains the materials are. It doesn't care what you do with these materials as long as it is all what you need.

We will try to describe this pizza making progress programmatically using C++ and Haskell.

Let's start with a complete C++ program
\begin{mdframed}[leftline=false,rightline=false,backgroundcolor=magenta!10,nobreak=false]
  \inputminted[linenos=true,breaklines,breaksymbolleft=,obeytabs=true,tabsize=2]{C}{../cpp/main.cpp}
\end{mdframed}

Output of this program
\begin{mdframed}[leftline=false,rightline=false,backgroundcolor=teal!10,nobreak=true]
  \inputminted[linenos=true,breaklines,breaksymbolleft=,obeytabs=true,tabsize=2]{text}{../cpp/main.txt}
\end{mdframed}

Nice! Let's do this again, but with Haskell
\begin{mdframed}[leftline=false,rightline=false,backgroundcolor=magenta!10,nobreak=false]
  \inputminted[linenos=true,breaklines,breaksymbolleft=,obeytabs=true,tabsize=2]{Haskell}{../hs/main.hs}
\end{mdframed}

Output of this program
\begin{mdframed}[leftline=false,rightline=false,backgroundcolor=teal!10,nobreak=true]
  \inputminted[linenos=true,breaklines,breaksymbolleft=,obeytabs=true,tabsize=2]{text}{../hs/main.txt}
\end{mdframed}

\newpage
\section{Conclusion}
With the pizza making out of the way, the two approaches feels like one.
Obviously, the process must be the same, but the way the pizza, or the data, here is treated differently.

With that said, here is an overkill comparison table.

\begin{center}
  \begin{tabularx}{\textwidth}{l*{2}{X}}
    \toprule
         & OOP                                                                                                                                                                                                    & FP                                                                                                                                                             \\
    \cmidrule(lr){2-3}
    Pros & OOP objects contains both the data (attributes) and things that it can do (methods)                                                                                                                    & FP decouples the data from the functions                                                                                                                       \\
         & Any changes that are applied is reflected on the object itself                                                                                                                                         & Data in FP is not intended to change, if it does, new data is just created                                                                                     \\
         & OOP organize everything into hierarchies of abstract objects                                                                                                                                           & FP is natural for the human brain, as in, our thought process is centered around ``doing'' things                                                              \\
    \midrule
    Cons & OOP objects are abstract and potentially complex, wasting a lot of time in abstractions and design patterns instead of solving the problem                                                             & A good FP program requires strict state and I/O management, which requires brain-breaking abstraction inversion                                                \\
         & OOP encourages sharing of mutable state and introduces additional complexity with its numerous design pattern, thus making common development practices, like refactoring and testing, needlessly hard & FP data objects are immutable, meaning that it is bad for simulating anything that involves objects that change like GUI                                       \\
         & Concurrency is basically impossible because the output of a method depends on the state of the object, unless we make the object immutable, which is FP                                                & FP has a strong mathematical foundation, meaning that it has a high barrier to entry, requiring you to learn concepts from abstract algebra to category theory \\
    \bottomrule
  \end{tabularx}
\end{center}

You can find the full source code of this report on \href{https://github.com/Smithienious/CO2039-Final}{my GitHub repository}.

\newpage
\section*{Easter egg}
Congratulations! You have actually read my report.
This section serves no more than empty space that I type when I can't get some thing out of my mind.

Yes, I am aware that this section creates the feeling that I am not taking this report seriously.
Yes, I am aware that the comparison table is not very related to the pizza problem, but FP is not something anyone can learn in 14 weeks, not to add that I had to learn it myself during the course.
The teacher was not very enthusiastic when it comes to FP either.

During the one-week duration to write this report, I was told that someone posted this report \href{https://www.chegg.com/homework-help/questions-and-answers/wikipedia-object-oriented-programming-template-method-one-behavioral-design-patterns-ident-q83199850}{on Chegg}.
Congratulations to whoever created this assignment, I guess.

With that having been said, please enjoy this little abomination I created in Python.
\begin{code}{Python}
class iostream:
    def __lshift__(self, other):
        print(other, end='')
        return self

    def __repr__(self):
        return ''


if __name__ == "__main__":
    cout = iostream()
    endl = '\n'
    cout << "Hello" << ", " << "World!" << endl
\end{code}

Thanks, and have fun.

\end{document}
